% GENERATED FILE. Edit canonical lessons, then run npm run generate:books.
% canonical-source-hash: fnv1a64:3378ca7369f3844f
% canonical-lessons: TA-C352-erpatuttu, TA-C352-takkam, TA-C352-karani, TA-C352-valivaku, TA-C352-pati

\chapter{Causes and Effects}
\label{ch:ta-causes-and-effects}
\hlchaptermodality{352}

\begin{chapteropening}
\textbf{By the end of this chapter:} \emph{I can say what caused something and what it led to.}

Complete the last lesson of chapter 352: I can say what caused something and what it led to.
\end{chapteropening}

\section[\texorpdfstring{\ta{ஏற்படுத்து} (ēṟpaṭuttu)}{ஏற்படுத்து (ēṟpaṭuttu)}]{\ta{ஏற்படுத்து} (ēṟpaṭuttu) — to cause; to bring about}

\label{lesson:TA-C352-erpatuttu}



\begin{quote}
Before the new one: say the Tamil for that is to say, then the Tamil for an example.
\end{quote}

\subsection*{You'll want to know: \ta{ஏற்படுத்து}}
\textbf{\ta{ஏற்படுத்து}} — \emph{ēṟpaṭuttu} — \textquotedblleft{}to cause; to bring about\textquotedblright{}.

\begin{quote}
\ta{அந்தச்} \ta{செய்தி} \ta{மக்களிடையே} \ta{குழப்பத்தை} \ta{ஏற்படுத்தியது}.
\end{quote}

\begin{quote}
\emph{antac ceyti makkaḷiṭaiyē kuḻappattai ēṟpaṭuttiyatu.}
\end{quote}

\begin{quote}
That news caused confusion among the people.
\end{quote}

Past (I …): \ta{ஏற்படுத்தினேன்}, \emph{ēṟpaṭuttiṉēṉ}.

The result takes -\ta{ஐ}: \ta{குழப்பத்தை} \ta{ஏற்படுத்து}. Its intransitive partner is \ta{ஏற்படு}, to come about.

\begin{cousinweb}
\textbf{\ta{ஏற்படுத்து}} (\emph{ēṟpaṭuttu}) is native Dravidian: DEDR 905 lists Tamil \emph{ēṟpaṭuttu}, \textquotedblleft{}to create, establish\textquotedblright{}, beside \emph{ēṟpaṭu}, \textquotedblleft{}to come into existence\textquotedblright{}, both from \emph{ēl}, \textquotedblleft{}to happen, to receive\textquotedblright{}.
\end{cousinweb}

\begin{grammarlens}[title={What you've built}]
One more word for giving reasons and opinions.
\end{grammarlens}

\subsection*{Your turn}
\begin{itemize}
\raggedright
  \item \emph{Say it:} \emph{ēṟpaṭuttu}
  \item \emph{Say it:} \emph{ēṟpaṭuttu}, once more
  \item \emph{Say it:} \emph{antac ceyti makkaḷiṭaiyē kuḻappattai ēṟpaṭuttiyatu.} — That news caused confusion among the people.
  \item \emph{Recall:} say the Tamil for that is to say, then the Tamil for an example, then say \emph{ēṟpaṭuttu} again
\end{itemize}

\begin{tcolorbox}[breakable,colback=teal!4,colframe=teal!35!black,title={Before you move on}]
What does \ta{ஏற்படுத்து} mean? (\textquotedblleft{}To cause; to bring about\textquotedblright{}.) Say it once more.
\end{tcolorbox}
\section[\texorpdfstring{\ta{தாக்கம்} (tākkam)}{தாக்கம் (tākkam)}]{\ta{தாக்கம்} (tākkam) — an impact; an influence}

\label{lesson:TA-C352-takkam}



\begin{quote}
Before the new one: say the Tamil for an example, then the Tamil for to cause.
\end{quote}

\subsection*{You'll want to know: \ta{தாக்கம்}}
\textbf{\ta{தாக்கம்}} — \emph{tākkam} — \textquotedblleft{}an impact; an influence\textquotedblright{}.

\begin{quote}
\ta{இந்தப்} \ta{புத்தகம்} \ta{என்} \ta{மீது} \ta{பெரிய} \ta{தாக்கத்தை} \ta{ஏற்படுத்தியது}.
\end{quote}

\begin{quote}
\emph{intap puttakam eṉ mītu periya tākkattai ēṟpaṭuttiyatu.}
\end{quote}

\begin{quote}
This book had a big impact on me.
\end{quote}

X \ta{மீது} \ta{தாக்கம்}, an impact on X.

\begin{cousinweb}
\textbf{\ta{தாக்கம்}} (\emph{tākkam}) is native Dravidian: DEDR 3150 lists Tamil \emph{tākku}, \textquotedblleft{}to collide, strike against\textquotedblright{}, and \emph{tākkam}, \textquotedblleft{}a striking\textquotedblright{}, the noun that now means \textquotedblleft{}impact\textquotedblright{}.
\end{cousinweb}

\begin{grammarlens}[title={What you've built}]
One more word for giving reasons and opinions.
\end{grammarlens}

\subsection*{Your turn}
\begin{itemize}
\raggedright
  \item \emph{Say it:} \emph{tākkam}
  \item \emph{Say it:} \emph{tākkam}, once more
  \item \emph{Say it:} \emph{intap puttakam eṉ mītu periya tākkattai ēṟpaṭuttiyatu.} — This book had a big impact on me.
  \item \emph{Recall:} say the Tamil for an example, then the Tamil for to cause, then say \emph{tākkam} again
\end{itemize}

\begin{tcolorbox}[breakable,colback=teal!4,colframe=teal!35!black,title={Before you move on}]
What does \ta{தாக்கம்} mean? (\textquotedblleft{}An impact; an influence\textquotedblright{}.) Say it once more.
\end{tcolorbox}
\section[\texorpdfstring{\ta{காரணி} (kāraṇi)}{காரணி (kāraṇi)}]{\ta{காரணி} (kāraṇi) — a factor}

\label{lesson:TA-C352-karani}



\begin{quote}
Before the new one: say the Tamil for to cause, then the Tamil for an impact.
\end{quote}

\subsection*{You'll want to know: \ta{காரணி}}
\textbf{\ta{காரணி}} — \emph{kāraṇi} — \textquotedblleft{}a factor\textquotedblright{}.

\begin{quote}
\ta{விலை} \ta{உயர்வுக்குப்} \ta{பல} \ta{காரணிகள்} \ta{இருக்கின்றன}.
\end{quote}

\begin{quote}
\emph{vilai uyarvukkup pala kāraṇikaḷ irukkiṉṟaṉa.}
\end{quote}

\begin{quote}
There are many factors behind the rise in prices.
\end{quote}

Plural: \ta{காரணிகள்}, \emph{kāraṇikaḷ}.

X-\ta{க்குக்} \ta{காரணி}, a factor in X. The plural is \ta{காரணிகள்}.

\begin{cousinweb}
\textbf{\ta{காரணி}} (\emph{kāraṇi}) is built on Sanskrit \emph{kāraṇa}, \textquotedblleft{}cause\textquotedblright{} (CDIAL 3057), the source of \ta{காரணம்}, with the ending -\ta{இ}: \textquotedblleft{}that which acts as a cause\textquotedblright{}, the word for a factor.
\end{cousinweb}

\begin{grammarlens}[title={What you've built}]
One more word for giving reasons and opinions.
\end{grammarlens}

\subsection*{Your turn}
\begin{itemize}
\raggedright
  \item \emph{Say it:} \emph{kāraṇi}
  \item \emph{Say it:} \emph{kāraṇi}, once more
  \item \emph{Say it:} \emph{vilai uyarvukkup pala kāraṇikaḷ irukkiṉṟaṉa.} — There are many factors behind the rise in prices.
  \item \emph{Recall:} say the Tamil for to cause, then the Tamil for an impact, then say \emph{kāraṇi} again
\end{itemize}

\begin{tcolorbox}[breakable,colback=teal!4,colframe=teal!35!black,title={Before you move on}]
What does \ta{காரணி} mean? (\textquotedblleft{}A factor\textquotedblright{}.) Say it once more.
\end{tcolorbox}
\section[\texorpdfstring{\ta{வழிவகு} (vaḻivaku)}{வழிவகு (vaḻivaku)}]{\ta{வழிவகு} (vaḻivaku) — to lead to; to pave the way for}

\label{lesson:TA-C352-valivaku}



\begin{quote}
Before the new one: say the Tamil for an impact, then the Tamil for a factor.
\end{quote}

\subsection*{You'll want to know: \ta{வழிவகு}}
\textbf{\ta{வழிவகு}} — \emph{vaḻivaku} — \textquotedblleft{}to lead to; to pave the way for\textquotedblright{}.

\begin{quote}
\ta{இந்த} \ta{முடிவு} \ta{பல} \ta{பிரச்சினைகளுக்கு} \ta{வழிவகுத்தது}.
\end{quote}

\begin{quote}
\emph{inta muṭivu pala piracciṉaikaḷukku vaḻivakuttatu.}
\end{quote}

\begin{quote}
This decision led to many problems.
\end{quote}

Past (it …): \ta{வழிவகுத்தது}, \emph{vaḻivakuttatu}.

The result takes -\ta{க்கு}: \ta{பிரச்சினைக்கு} \ta{வழிவகு}.

\begin{cousinweb}
\textbf{\ta{வழிவகு}} (\emph{vaḻivaku}) is native: \emph{vaḻi}, \textquotedblleft{}way, path\textquotedblright{} (DEDR 5297), and \emph{vaku}, \textquotedblleft{}to separate, apportion, lay out\textquotedblright{} (DEDR 5202): to lay out the way.
\end{cousinweb}

\begin{grammarlens}[title={What you've built}]
One more word for giving reasons and opinions.
\end{grammarlens}

\subsection*{Your turn}
\begin{itemize}
\raggedright
  \item \emph{Say it:} \emph{vaḻivaku}
  \item \emph{Say it:} \emph{vaḻivaku}, once more
  \item \emph{Say it:} \emph{inta muṭivu pala piracciṉaikaḷukku vaḻivakuttatu.} — This decision led to many problems.
  \item \emph{Recall:} say the Tamil for an impact, then the Tamil for a factor, then say \emph{vaḻivaku} again
\end{itemize}

\begin{tcolorbox}[breakable,colback=teal!4,colframe=teal!35!black,title={Before you move on}]
What does \ta{வழிவகு} mean? (\textquotedblleft{}To lead to; to pave the way for\textquotedblright{}.) Say it once more.
\end{tcolorbox}
\section[\texorpdfstring{\ta{பாதி} (pāti)}{பாதி (pāti)}]{\ta{பாதி} (pāti) — to affect; to harm}

\label{lesson:TA-C352-pati}



\begin{quote}
Before the new one: say the Tamil for a factor, then the Tamil for to lead to.
\end{quote}

\subsection*{You'll want to know: \ta{பாதி}}
\textbf{\ta{பாதி}} — \emph{pāti} — \textquotedblleft{}to affect; to harm\textquotedblright{}.

\begin{quote}
\ta{வெள்ளம்} \ta{பல} \ta{விவசாயிகளைப்} \ta{பாதித்தது}.
\end{quote}

\begin{quote}
\emph{veḷḷam pala vivacāyikaḷaip pātittatu.}
\end{quote}

\begin{quote}
The flood affected many farmers.
\end{quote}

Past (it …): \ta{பாதித்தது}, \emph{pātittatu}.

Who is affected takes -\ta{ஐ}. Not the noun \ta{பாதி}, a half: that is another word.

\begin{cousinweb}
\textbf{\ta{பாதி}} (\emph{pāti}), \textquotedblleft{}to affect\textquotedblright{}, is a verb built on the Sanskrit root \emph{bādh-}, \textquotedblleft{}to press, oppress, harass\textquotedblright{} (Monier-Williams). Tamil writes the Sanskrit b as \ta{ப}.
\end{cousinweb}

\begin{grammarlens}[title={What you've built}]
One more word for giving reasons and opinions.
\end{grammarlens}

\subsection*{Your turn}
\begin{itemize}
\raggedright
  \item \emph{Say it:} \emph{pāti}
  \item \emph{Say it:} \emph{pāti}, once more
  \item \emph{Say it:} \emph{veḷḷam pala vivacāyikaḷaip pātittatu.} — The flood affected many farmers.
  \item \emph{Recall:} say the Tamil for a factor, then the Tamil for to lead to, then say \emph{pāti} again
\end{itemize}

\begin{tcolorbox}[breakable,colback=teal!4,colframe=teal!35!black,title={Before you move on}]
What does \ta{பாதி} mean? (\textquotedblleft{}To affect; to harm\textquotedblright{}.) Say it once more.
\end{tcolorbox}
